% TODO make hw stylesheet
\documentclass[11pt]{scrartcl}
\usepackage[a4paper, total={6in, 8in}]{geometry}
\usepackage[usenames,dvipsnames,svgnames]{xcolor}
\usepackage[shortlabels]{enumitem}
\usepackage[framemethod=TikZ]{mdframed}
\usepackage{amsmath,amssymb,amsthm}
\usepackage[colorlinks]{hyperref}
\usepackage{multicol}
\usepackage{tabularx}

\hypersetup{
	colorlinks=true,
	linkcolor=blue,
	filecolor=magenta,      
	urlcolor=magenta,
	pdftitle={MAT 322 HW}
}

\usepackage{microtype}
\usepackage{mathtools}
\usepackage[headsepline]{scrlayer-scrpage}
\usepackage{thmtools}
\usepackage[textsize=footnotesize]{todonotes}

\mdfdefinestyle{mdbluebox}{roundcorner=10pt,innerbottommargin=9pt,
	linecolor=blue,backgroundcolor=TealBlue!5,}
\declaretheoremstyle[headfont=\sffamily\bfseries\color{MidnightBlue},
mdframed={style=mdbluebox},]{thmbluebox}

\mdfdefinestyle{mdbloobox}{frametitlefont=\bfseries,innerbottommargin=8pt,
	nobreak=true,backgroundcolor=TealBlue!5,linecolor=blue,}
\declaretheoremstyle[headfont=\bfseries\color{MidnightBlue},
mdframed={style=mdbloobox},headpunct={\\[3pt]},postheadspace=0pt,]{thmbloobox}

\mdfdefinestyle{mdredbox}{frametitlefont=\bfseries,innerbottommargin=8pt,
	nobreak=true,backgroundcolor=Salmon!5,linecolor=RawSienna,}
\declaretheoremstyle[headfont=\bfseries\color{RawSienna},
mdframed={style=mdredbox},headpunct={\\[3pt]},postheadspace=0pt,]{thmredbox}

\mdfdefinestyle{mdgreenbox}{linecolor=ForestGreen,backgroundcolor=ForestGreen!5,
	linewidth=2pt,rightline=false,leftline=true,topline=false,bottomline=false,}
\declaretheoremstyle[headfont=\bfseries\sffamily\color{ForestGreen!70!black},
mdframed={style=mdgreenbox},headpunct={ --- },]{thmgreenbox}

\mdfdefinestyle{mdorangebox}{linecolor=RedOrange,backgroundcolor=RedOrange!5,
	linewidth=2pt,rightline=false,leftline=true,topline=false,bottomline=false,}
\declaretheoremstyle[headfont=\bfseries\sffamily\color{RedOrange!70!black},
mdframed={style=mdorangebox},headpunct={ --- },]{thmorangebox}

\mdfdefinestyle{mdblackbox}{linecolor=black,backgroundcolor=RedViolet!5!gray!5,
	linewidth=3pt,rightline=false,leftline=true,topline=false,bottomline=false,}
\declaretheoremstyle[mdframed={style=mdblackbox}]{thmblackbox}

\declaretheoremstyle[spaceabove=6pt,spacebelow=6pt]{basehead}

\declaretheorem[style=basehead,name=Problem]{problem}
\declaretheorem[style=thmredbox,name=Problem,numbered=no]{problem*}
\declaretheorem[style=thmbloobox,name=Setup]{setup} % for config geo
\declaretheorem[style=thmredbox,name=Example,sibling=problem]{example}
\declaretheorem[style=thmredbox,name=$\bigstar$ Problem,sibling=problem]{niceprob}
\declaretheorem[style=thmbluebox,name=Theorem,sibling=problem]{theorem}
\declaretheorem[style=thmbluebox,name=Corollary,sibling=theorem]{corollary}
\declaretheorem[style=thmbluebox,name=Proposition,sibling=theorem]{prop}
\declaretheorem[style=thmbluebox,name=Lemma,numberwithin=problem]{lemma}
\declaretheorem[style=thmbluebox,name=Theorem,numbered=no]{theorem*}
\declaretheorem[style=thmbluebox,name=Lemma,numbered=no]{lemma*}
\declaretheorem[style=thmgreenbox,name=Claim,numberwithin=problem]{claim}
\declaretheorem[style=thmgreenbox,name=Claim,numbered=no]{claim*}
\declaretheorem[style=thmblackbox,name=Remark,numberwithin=problem]{remark}
\declaretheorem[style=thmblackbox,name=Remark,numbered=no]{remark*}
\declaretheorem[style=thmgreenbox,name=Definition,sibling=theorem]{definition}
\declaretheorem[style=thmgreenbox,name=Definition,numbered=no]{definition*}
\declaretheorem[style=thmorangebox,name=Warning,sibling=theorem]{warning}
\declaretheorem[style=thmorangebox,name=Warning,numbered=no]{warning*}
\declaretheorem[style=thmorangebox,name=Note,numbered=no]{note}
\declaretheorem[style=thmblackbox,name=Exercise,sibling=problem]{exercise}
\declaretheorem[style=thmblackbox,name=Exercise,numbered=no]{exercise*}
\declaretheorem[style=thmblackbox,name=Question,sibling=problem]{question}
\declaretheorem[style=thmblackbox,name=Question,numbered=no]{question*}

\addtolength{\textheight}{3.14cm}
\providecommand{\ol}{\overline}
\providecommand{\ul}{\underline}
\providecommand{\eps}{\varepsilon}
\providecommand{\half}{\frac{1}{2}}
\providecommand{\dang}{\measuredangle} %% Directed angle
\providecommand{\CC}{\mathbb C}
\providecommand{\FF}{\mathbb F}
\providecommand{\NN}{\mathbb N}
\providecommand{\QQ}{\mathbb Q}
\providecommand{\RR}{\mathbb R}
\providecommand{\ZZ}{\mathbb Z}
\providecommand{\dg}{^\circ}
\providecommand{\ii}{\item}
\providecommand{\setdiff}{\smallsetminus}
\providecommand{\alert}[1]{{\sffamily\textbf{\textcolor{blue}{#1}}}}
\providecommand{\inv}{^{-1}}
\providecommand{\mb}[1]{\mathbf{#1}}
\providecommand{\dif}{\;d}

\reversemarginpar

\newenvironment{soln}{\begin{proof}[Solution]}{\end{proof}}
\newenvironment{walkthrough}{\noindent \textbf{Walkthrough.}}{\medskip}
\newlist{walk}{enumerate}{3}
\setlist[walk]{label=\bfseries (\alph*)}

\providecommand{\pow}{\operatorname{Pow}}
\providecommand{\vocab}[1]{{\textbf{\textcolor{ForestGreen}{#1}}}}
\providecommand{\lcm}{\operatorname{lcm}}
\providecommand{\abs}[1]{\left|#1\right|}

\usepackage{asymptote}
\begin{asydef}
	defaultpen(fontsize(10pt));
	unitsize(1cm);
	size(8cm); // set a reasonable default
	usepackage("amsmath");
	usepackage("amssymb");
	settings.tex="pdflatex";
	settings.outformat="pdf";
	// Replacement for olympiad+cse5 which is not standard
	import geometry;
	// recalibrate fill and filldraw for conics
	void filldraw(picture pic = currentpicture, conic g, pen fillpen=defaultpen, pen drawpen=defaultpen)
	{ filldraw(pic, (path) g, fillpen, drawpen); }
	void fill(picture pic = currentpicture, conic g, pen p=defaultpen)
	{ filldraw(pic, (path) g, p); }
	// some geometry
	pair foot(pair P, pair A, pair B) { return foot(triangle(A,B,P).VC); }
	pair orthocenter(pair A, pair B, pair C) { return orthocentercenter(A,B,C); }
	pair centroid(pair A, pair B, pair C) { return (A+B+C)/3; }
	// cse5 abbreviations
	path CP(pair P, pair A) { return circle(P, abs(A-P)); }
	path CR(pair P, real r) { return circle(P, r); }
	pair IP(path p, path q) { return intersectionpoints(p,q)[0]; }
	pair OP(path p, path q) { return intersectionpoints(p,q)[1]; }
	path Line(pair A, pair B, real a=0.6, real b=a) { return (a*(A-B)+A)--(b*(B-A)+B); }
	// cse5 more useful functions
	picture CC() {
		picture p=rotate(0)*currentpicture;
		currentpicture.erase();
		return p;
	}
	pair MP(Label s, pair A, pair B = plain.S, pen p = defaultpen) {
		Label L = s;
		L.s = "$"+s.s+"$";
		label(L, A, B, p);
		return A;
	}
	pair Drawing(Label s = "", pair A, pair B = plain.S, pen p = defaultpen) {
		dot(MP(s, A, B, p), p);
		return A;
	}
	path Drawing(path g, pen p = defaultpen, arrowbar ar = None) {
		draw(g, p, ar);
		return g;
	}
\end{asydef}

\usepackage{mathrsfs}

\ihead{\footnotesize MAT 322 HW01}
\ohead{\footnotesize Last compiled \today}

\title{\vspace{-2em}MAT 322 HW01}
\author{\large Aditya Pahuja}
\date{\large Due 02/04}
\begin{document}
\maketitle

\begin{problem}
	Prove the \emph{Cauchy-Schwarz} inequality:
	for any two vectors $x,y\in\RR^n$,
	\[|\langle x,y\rangle| \le \|x\|\|y\|.\]
\end{problem}

\begin{soln}
	Since both sides of the inequality are nonnegative,
	squaring gives the equivalent inequality
	\[|\langle x,y\rangle|^2 \le
	\|x\|^2\|y\|^2.\]
	Expanding in terms of the components tells us that we need to prove
	\[(x_1y_1 + x_2y_2 + \cdots + x_ny_n)^2 \le
	(x_1^2 + x_2^2 + \cdots + x_n^2)(y_1^2 + y_2^2 + \cdots + y_n^2)\]
	for any $x_i,y_i\in\RR$, $i=1,2,\dots,n$.
	The left-hand side is equal to
	\[\sum_{i=1}^n (x_iy_i)^2 + 2\sum_{1\le i<j\le n} x_iy_i\cdot x_jy_j,\]
	while the right-hand side is equal to
	\[\sum_{i=1}^n x_i^2 y_i^2 + \sum_{1\le i<j\le n} x_i^2y_j^2 + x_j^2y_i^2.\]
	Noting that the first addends in the two expressions are the same,
	the inequality reduces to proving that
	\[2\sum_{1\le i<j\le n} x_iy_i\cdot x_jy_j\le
	\sum_{1\le i<j\le n} x_i^2y_j^2 + x_j^2y_i^2.\]
	Indeed, for any real numbers $a$, $b$, $c$, $d$, we have
	\[a^2b^2 + c^2d^2 \ge 2abcd
	\iff a^2b^2 - 2abcd + c^2d^2 = (ab-cd)^2 \ge 0,\]
	which is true because squares of real numbers are nonnegative
	(alternatively, this is just two-variable AM-GM).
	Thus, taking $(a,b,c,d) = (x_i,y_j,x_j,y_i)$ gives us the inequality
	\[x_i^2y_j^2 + x_j^2y_i^2 \ge 2x_iy_jx_jy_i.\]
	Summing over all pairs $(i,j)$ with $i$, $j$ integers satisfying
	$1\le i<j\le n$ therefore yields the inequality
	\[\sum_{1\le i<j\le n} x_i^2y_j^2 + x_j^2y_i^2 \ge
	\sum_{1\le i<j\le n} 2x_iy_jx_jy_i =
	2\sum_{1\le i<j\le n} x_iy_i\cdot x_jy_j,\]
	which is what we wanted.
\end{soln}

\pagebreak
\begin{problem}
	Prove the following statements:
	\begin{enumerate}[(a)]
		\ii For any $r>0$, the $r$-ball around $x$, defined as
		\[B_r(x) = \{y\in\RR^n : \|x-y\| < r\},\]
		is an open set.
		\ii A function $f\colon\RR^n\to\RR^m$ is continuous at $x\in\RR^n$
		if and only if
		\[\lim_{y\to x} f(y) = f(x).\]
		\ii If $f\colon\RR^n\to\RR$ is continuous at $x\in\RR^n$
		and $f(x)\ne0$, then $1/f$ is a continuous function at $x$.
		\ii If $f\colon\RR^n\to\RR^m$ is continuous at $x\in\RR$,
		then its norm $\|f\|\colon\RR^n\to\RR$
		is a continuous function at $x$. Is the converse true?
	\end{enumerate}
\end{problem}

\begin{soln}
	(a) We wish to show that for any $p\in B_r(x)$, there is
	an open ball centered at $p$ that is a subset of $B_r(x)$.
	Consider the ball $B := B_{r-\|x-p\|}(p)$. If $y\in B$, then
	\[\|x-y\| \le \|x-p\| + \|p-y\| < \|x-p\| + (r - \|x-p\|) = r,\]
	where the first inequality is the triangle inequality
	and the second follows from the definition of open ball, since
	$y$ being in the ball $B$ means that its distance to
	the ball's center $p$ is less than the radius $r - \|x-p\|$.
	That is, the distance between $x$ and $y$ is smaller than $r$,
	so $y\in B$ implies $y\in B_r(x)$ as well, so $B_r(x)$ is open.\\
	
	(b) The statement
	\[\lim_{y\to x}f(y) = f(x)\]
	means that for every $\eps > 0$,
	there exists a $\delta > 0$ such that $0 < \|x-y\| < \delta$
	implies $\|f(x) - f(y)\| < \eps$. This is the exact same thing
	as continuity at $x$; the edge case $\|x-y\| = 0$ is irrelevant
	because obviously $\|f(x) - f(y)\|$ would also be $0 < \eps$.\\
	
	(c) Since $f$ is continuous at $x$, there exists a $\alpha$ such that
	\[|x-y|<\alpha\implies |f(x)-f(y)| < \frac{|f(x)|}{2}.\]
	This in particular means that $|f(y)| \ge \frac{|f(x)|}{2}$
	when $|x-y|<\alpha$, since otherwise
	\[|f(x)| \le |f(x) - f(y)| + |f(y)| <
	\frac{|f(x)|}{2} + \frac{|f(x)|}{2} = |f(x)|.\]
	
	Define $C = \frac{f(x)^2}{2}$ and let $\eps > 0$.
	Since $f$ is continuous at $x$, there is a $0 < \delta'$ such that
	$|f(x)-f(y)| < C\eps$ whenever $|x-y|<\delta'$.
	Define $\delta = \min(\alpha,\delta')$. Then,
	\[|x-y| < \delta\implies |f(x) - f(y)| < C\eps \le |f(x)f(y)|\eps,\]
	and the latter inequality rewrites to
	\[\left|\frac{1}{f(y)} - \frac{1}{f(x)}\right| < \eps\]
	upon division by $|f(x)f(y)|$, which proves continuity.
	\\
	
	(d) Given $\eps > 0$, since $f$ is continuous at $x$,
	let $\delta > 0$ be such that
	\[\|x-y\| < \delta \implies \|f(x) - f(y)\| < \eps.\]
	We want to prove that there exists a $\delta'>0$ such that
	\[\|x-y\| < \delta'\implies \| \|f(x)\| - \|f(y)\|\|
	= \abs{\|f(x)\| - \|f(y)\|} < \eps.\]
	The triangle inequality says that
	\[\|f(x)\| \le \|f(y)\| + \|f(x) - f(y)\|,\qquad
	\|f(y)\| \le \|f(x)\| + \|f(y) - f(x)\|,\]
	so we get
	\[-\|f(x) - f(y)\| \le \|f(x)\| - \|f(y)\| \le \|f(x) - f(y)\|\]
	which is equivalent to
	\[\abs{\|f(x)\| - \|f(y)\|} \le \|f(x) - f(y)\|.\]
	This means we can actually just take $\delta' := \delta$
	to find that
	\[\|x-y\|<\delta \implies \|\|f(x)\| - \|f(y)\|\|
	\le \|f(x)-f(y)\| < \eps.\]
	This proves that $\|f\|$ is continuous at $x$,
	as, for any $\eps>0$, we have found a $\delta'>0$
	such that $y\in B_{\delta'}(x)$ implies $f(y)\in B_\eps(x)$.
	
	The converse statement, however, is false.
	For example, take $f\colon\RR\to\RR$ to be the function
	$f(x) = 1$ for $x\ge 0$ and $f(x) = -1$ for $x < 0$.
	The norm is the constant function that is always equal to $1$,
	which is obviously continuous, but $f$ isn't continuous
	because the interval $(0,\infty)$ is an open set
	whose pre-image $f\inv((0,\infty)) = [0,\infty)$ is not open.
\end{soln}

\begin{problem}
	Let $f\colon\RR\to\RR$ be defined by
	\[f(x) = \begin{cases}
		\sin(x) & \text{if }x\in\QQ\\
		0 & \text{if }x\in\RR\setdiff\QQ.
	\end{cases}\]
	At which points in $\RR$ is $f$ continuous?
\end{problem}

\begin{soln}
	The function is continuous at \framebox{integer multiples of $\pi$},
	i.e. the points at which $\sin(x) = 0$.
	To check that $f$ is continuous at $x = n\pi$ for any $n\in\ZZ$,
	let $\eps > 0$. Since $\sin$ is continuous on all of $\RR$,
	there exists a $\delta > 0$ such that
	\[|n\pi-y| < \delta\implies |\sin(n\pi) - \sin(y)|
	= |\sin (y)| < \eps.\]
	Then, $|\sin(y)| = \max\{|\sin (y)|, 0\} \ge |f(y)|$
	(for $f(y)$ is either equal to $0$ or $\sin(y)$)
	and $f(n\pi) = 0$ on account of $n\pi$ being irrational, meaning
	\[|n\pi - y| < \delta \implies |f(n\pi) - f(y)|
	= |f(y)| \le |\sin(y)| < \eps.\]
	We have therefore shown that the definition of continuity at $x$
	holds when $x = n\pi$ for an integer $n$.
	
	Now we have to show that $f$ is discontinuous at all other points.
	If $x$ is not an integer multiple of $\pi$, then $\sin(x)\ne 0$,
	so pick $\eps = |\sin(x)|/2>0$ and let $\delta > 0$ be arbitrary.
	Additionally, let $\delta' > 0$ be chosen so that
	$|x - y| < \delta'$ implies that $|\sin(x) - \sin(y)| < \eps/2$,
	which we can do by continuity of $\sin$.
	Since $\QQ$ and its complement are both dense in $\RR$,
	we can pick a rational $q$ and an irrational $r$
	that lie inside the interval $(x-m,x+m)$,
	where $m$ is the smaller of $\delta$, $\delta'$.
	Then, either $f(x) = 0$, in which case $|x - r| < m$ implies
	\[|f(x) - f(r)| = |\sin(r)| > |\sin(x)| - |\sin(x) - \sin(r)|
	> 2\eps - \eps/2 > \eps,\]
	so $f$ can't be continuous at $x$,
	or $f(x) = \sin(x)$, in which case $|x - q| < m$ but
	\[|f(x) - f(q)| = |f(x)| = |\sin(x)| = 2\eps > \eps.\]
	In either case we have found a point in the interval $(x-\delta,x+\delta)$
	whose image under $f$ is not within $(f(x)-\eps,f(x)+\eps)$,
	so $f$ cannot be continuous at any $x$ outside $\pi\ZZ$.
\end{soln}

\begin{problem}
	Consider the function $f\colon\RR^2\to\RR$ given by $f(0,0)=0$ and
	\[f(x,y) = \frac{xy}{x^2 + y^4}\text{ if }(x,y)\ne(0,0).\]
	At which points in $\RR^2$ is $f$ continuous?
\end{problem}

\begin{soln}
	$f$ is continuous on $\boxed{\RR\setdiff\{(0,0)\}}$.
	
	First, we show that $f$ is continuous on $\RR\setdiff\{(0,0)\}$.
	This is pretty simple: the $\RR^2\to\RR$
	functions $a(x,y) = x$ and $b(x,y) = y$
	are very obviously continuous on $\RR$ (if $U\subseteq\RR$ is open,
	then $a\inv(U) = U\times\RR$ and $b\inv(U) = \RR\times U$ are open),
	and then we can use the fact that sums, products, and quotients
	of continuous functions are continuous (with the caveat that
	the divisor in a quotient has to be nonzero) to conclude that
	\[f(x,y) = \frac{xy}{x^2 + y^4}\]
	is continuous whenever $x^2+y^4\ne 0$, i.e. all of $\RR\setdiff\{(0,0)\}$.
	
	To show that $f$ is discontinuous at $(0,0)$, consider the sequence
	$a_n = (\frac1n, \frac1n)$ for $n\in\NN$.
	This sequence has limit $(0,0)$ in $\RR^2$, but
	\[\lim_{n\to\infty} f(a_n) = \lim_{n\to\infty} \frac{1/n^2}{1/n^2 + 1/n^4}
	= \lim_{n\to\infty} \frac{n^2}{n^2 + 1} = 1\ne f(0,0)\]
	which cannot hold if $f$ is continuous at $(0,0)$.
\end{soln}


























\end{document}